\begin{helper}[Natural reduced-vertex edge-mass exponential remainder]
For every constant \(c>0\), and every natural reduced-vertex count
\(m(n)=\text{\(\noderef{TwoBiteNaturalReducedVertexCount}\)}(n)\), one has
\[
  m(n)\exp\!\left(
    -c\,\text{\(\noderef{TwoBiteEdgeProbability}\)}(1/2,n)\,n
  \right)\longrightarrow 0 .
\]
\end{helper}
\begin{proof}
By \(\noderef{NaturalDegreeMassDominatesLogCube}\), applied with \(c^{-1}\),
eventually
\[
  c^{-1}(\log n)^3
  \le
  \text{\(\noderef{TwoBiteEdgeProbability}\)}(1/2,n)
  \text{\(\noderef{TwoBiteNaturalReducedVertexCount}\)}(n).
\]
Multiplying by \(c\) gives
\[
  (\log n)^3
  \le
  c\,\text{\(\noderef{TwoBiteEdgeProbability}\)}(1/2,n)\,
  \text{\(\noderef{TwoBiteNaturalReducedVertexCount}\)}(n).
\]
The floor definition of
\(\text{\(\noderef{TwoBiteNaturalReducedVertexCount}\)}(n)\), together with
the reduced-vertex formula from \(\noderef{TwoBiteReducedVertexCount}\) and
\(\noderef{TwoBiteStretch}\), gives
\(\text{\(\noderef{TwoBiteNaturalReducedVertexCount}\)}(n)\le n\) for all
large \(n\).  Since
\(\text{\(\noderef{TwoBiteEdgeProbability}\)}(1/2,n)\ge0\), the preceding
display implies eventually
\[
  (\log n)^3
  \le c\,\text{\(\noderef{TwoBiteEdgeProbability}\)}(1/2,n)\,n .
\]
Also \(m(n)\le n\).  Hence the displayed sequence is eventually bounded above
by
\[
  n\exp(-(\log n)^3)
  =
  \exp(\log n-(\log n)^3).
\]
The function \(y\mapsto y-y^3\) tends to \(-\infty\) as \(y\to\infty\), and
\(\log n\to\infty\), so the last expression tends to \(0\).  The original
sequence is nonnegative, and the squeeze theorem proves the claim.
\end{proof}
